Here is your exercise **cleanly formatted as a Markdown question sheet**, with **all text preserved** and structured for clarity.

---

# **Repeat 2007 — Exercise**

## **4. Weight‑loss and Two‑Sample Inference**

### **(a)**  
The weight of 8 people before and after a diet are given below (in kg):

| Person | 1 | 2 | 3 | 4 | 5 | 6 | 7 | 8 |
|--------|---|---|---|---|---|---|---|---|
| **Before** | 87 | 95 | 77 | 97 | 83 | 102 | 85 | 79 |
| **After**  | 85 | 91 | 76 | 93 | 83 | 99 | 84 | 78 |

i) **Test whether the diet reduces weight by 4 kg on average** at a significance level of 5%.

ii) **State the assumption made when carrying out this test**.

*(10 marks)*

---

### **(b)**  
The weight of 50 Americans and 100 Spaniards was observed.  
- Americans: mean = 84 kg, standard deviation = 15 kg  
- Spaniards: mean = 78 kg, standard deviation = 12 kg  

i) **By calculating the realisation of the appropriate test statistic**, determine whether it can be stated at a significance level of 1% that the mean weight of all Americans differs from the mean weight of all Spaniards.  
Clearly state your hypotheses.
